% Options for packages loaded elsewhere
\PassOptionsToPackage{unicode}{hyperref}
\PassOptionsToPackage{hyphens}{url}
\documentclass[
  a4paper,
]{ltjsarticle}
\usepackage{xcolor}
\usepackage[margin=25mm]{geometry}
\usepackage{amsmath,amssymb}
\setcounter{secnumdepth}{-\maxdimen} % remove section numbering
\usepackage{iftex}
\ifPDFTeX
  \usepackage[T1]{fontenc}
  \usepackage[utf8]{inputenc}
  \usepackage{textcomp} % provide euro and other symbols
\else % if luatex or xetex
  \usepackage{unicode-math} % this also loads fontspec
  \defaultfontfeatures{Scale=MatchLowercase}
  \defaultfontfeatures[\rmfamily]{Ligatures=TeX,Scale=1}
\fi
\usepackage{lmodern}
\ifPDFTeX\else
  % xetex/luatex font selection
\fi
% Use upquote if available, for straight quotes in verbatim environments
\IfFileExists{upquote.sty}{\usepackage{upquote}}{}
\IfFileExists{microtype.sty}{% use microtype if available
  \usepackage[]{microtype}
  \UseMicrotypeSet[protrusion]{basicmath} % disable protrusion for tt fonts
}{}
\makeatletter
\@ifundefined{KOMAClassName}{% if non-KOMA class
  \IfFileExists{parskip.sty}{%
    \usepackage{parskip}
  }{% else
    \setlength{\parindent}{0pt}
    \setlength{\parskip}{6pt plus 2pt minus 1pt}}
}{% if KOMA class
  \KOMAoptions{parskip=half}}
\makeatother
\setlength{\emergencystretch}{3em} % prevent overfull lines
\providecommand{\tightlist}{%
  \setlength{\itemsep}{0pt}\setlength{\parskip}{0pt}}
\usepackage{bookmark}
\IfFileExists{xurl.sty}{\usepackage{xurl}}{} % add URL line breaks if available
\urlstyle{same}
\hypersetup{
  pdftitle={無名複素定常波真空からの局在凝縮体・重心・質量的計量の読出し ――Fourier--Mellin 型匿名読出しによる凝縮強度、有限広がり、および真空状態ゲージ歪みの構成},
  pdfauthor={木原 範昭（Independent Researcher）　ORCID 0009-0004-6753-4020},
  hidelinks,
  pdfcreator={LaTeX via pandoc}}

\title{無名複素定常波真空からの局在凝縮体・重心・質量的計量の読出し
――Fourier--Mellin
型匿名読出しによる凝縮強度、有限広がり、および真空状態ゲージ歪みの構成}
\author{木原 範昭（Independent Researcher）　ORCID 0009-0004-6753-4020}
\date{2026年9月3日　v3.0　Version DOI 10.5281/zenodo.22282690}

\begin{document}
\maketitle

\subsection{要旨}\label{ux8981ux65e8}

本稿は、局在した粒子、質量、重心、背景時空を最初から仮定せず、無名な複素定常波の集合と複素二乗ゼロ閉包

\[
\sum_n z_n^2=0
\]

を出発点として、均質真空状態場の内部に局在した凝縮体が形成されたとき、その凝縮体の相対重心、有限広がり、凝縮強度、および周囲の真空状態ゲージの歪みを、外部座標や粒子ラベルを導入せずに読出す方法を構成する。

前報 {[}K1{]}
では、無名複素ゼロ閉包系から相対尺度、位相周期、対数的空間・時間読出し、および写像後の
\[\lambda\nu\]
型不変量を構成した。本稿では、その均質真空状態場を「複数複素定常波の海」として受け取り、そこに均質背景から識別可能な局在凝縮が生じた場合を扱う。

複素数

\[
z_n=r_n e^{i\theta_n}
\]

に対し、対数振幅

\[
X_n=\ln r_n
\]

と位相 \[\theta_n\] を用いて、無名集合の経験測度に対する Fourier--Mellin
型変換

\[
\mathcal A(s,m)
=
\frac{1}{N}
\sum_{n=1}^{N}
\exp\left[
-i\left(
s\ln r_n+m\theta_n
\right)
\right]
\]

を定義する。この変換は波の並べ替えに不変であり、共通振幅スケールおよび共通位相回転は変換全体の位相因子としてのみ現れるため、その強度

\[
P(s,m)=|\mathcal A(s,m)|^2
\]

は匿名性、共通スケール、および共通位相に不変である。

局在凝縮体の正規化変換を \[\chi_C\] とすると、その対数は変換原点近傍で

\[
\ln\chi_C(\mathbf q)
=
-i\mathbf q\cdot\boldsymbol\mu_C
-\frac12
\mathbf q^{\mathsf T}\Sigma_C\mathbf q
+O(|\mathbf q|^3)
\]

と展開できる。ここで \[\boldsymbol\mu_C\] は凝縮体の重心、\[\Sigma_C\]
は有限広がりと形状を表す共分散行列である。したがって重心はスペクトル位相の一次微分、広がりはスペクトル振幅の二次微分から読出せる。ただし絶対重心は共通ゲージ変換により消去され、物理的に読出せるのは別の凝縮体または均質真空背景に対する相対重心である。

さらに、凝縮前の均質真空スペクトル \[\mathcal A_0\] と凝縮後のスペクトル
\[\mathcal A\] の差

\[
\delta\mathcal A
=
\mathcal A-\mathcal A_0
\]

を導入すると、凝縮体の存在は真空状態場の局所的な非均質として定義できる。逆変換された
\[\delta\rho\]
は凝縮中心の外側にも広がり得るため、局在凝縮が周囲の真空状態ゲージを変形する条件を直接検証できる。本稿では、質量を外部パラメータとして置かず、凝縮強度、安定した相対重心、有限二次モーメント、および背景ゲージ変形を同一の状態変化から得る「質量的計量」を提案する。

補遺Aではさらに、前報の光的特殊解を \[P=P_0L^\alpha\]
へ一般化し、全殻対で \[\lambda\nu\]
が一定となるための必要十分条件がこの冪関係であることを示す。読出し単位を一致させると、\[\alpha<1\]、\[\alpha=1\]、\[\alpha>1\]
がそれぞれ spacelike、null、timelike
に対応する。また、整数巡回・各殻の厳密二乗ゼロ閉包・全体絶対収束を同時に満たす有理指数殻族を構成する。ただし、局在凝縮が実際に
\[\alpha_{\mathrm{eff}}>1\]
を生成するという接続は未導出の仮説として分離する。

本稿の主張は、現実の粒子質量や Einstein
方程式を既に導出したというものではない。示すのは、二乗ゼロ閉包した無名複素定常波真空に対して、局在凝縮体を匿名的に検出し、その相対重心、有限広がり、凝縮強度、および背景状態場の歪みを同一の調和解析から定義できる明示的な読出し枠組みである。

\textbf{キーワード}：無名複素数、二乗ゼロ閉包、定常波真空、局在凝縮、重心、質量、Fourier--Mellin
変換、特性関数、ゲージ不変読出し、創発物質

\begin{center}\rule{0.5\linewidth}{0.5pt}\end{center}

\subsection{1.
問題設定――均質真空から「物体」をどう読むか}\label{ux554fux984cux8a2dux5b9aux5747ux8ceaux771fux7a7aux304bux3089ux7269ux4f53ux3092ux3069ux3046ux8aadux3080ux304b}

前報 {[}K1{]}
では、空間、時間、波長、振動数を先に置かず、無名な複素数集合

\[
\mathcal Z
=
\{z_1,z_2,\ldots\}/S
\]

と二乗ゼロ閉包

\[
Q(\mathcal Z)
=
\sum_n z_n^2
=
0
\tag{1}
\]

から、相対的に何が読出し可能かを調べた。

式 (1) は任意の共通複素スケール

\[
z_n\mapsto \gamma z_n,
\qquad
\gamma\in\mathbb C^\times
\]

に対して

\[
Q\mapsto \gamma^2Q=0
\]

を保つ。したがって閉包条件そのものは絶対振幅スケールも絶対位相も選ばない。

本稿では、このような状態を多数の複素定常波からなる均質真空状態場として読む。重要なのは、これは「時空の中に存在する真空場」ではなく、読出し時空以前の状態集合であることである。

均質真空では、どの部分にも特別な中心を置かない。したがって、真空そのものには物体固有の重心は存在しない。

本稿の問いは次である。

\begin{quote}
無名な複素定常波の海の一部が凝縮し、背景から区別可能な局在構造を形成したとき、外部座標、粒子番号、既知の質量を与えずに、その凝縮体の重心、広がり、強度、および周囲の背景場の歪みを読出せるか。
\end{quote}

必要条件を先に整理すると、物質的凝縮体の候補は少なくとも以下を満たす必要がある。

\begin{enumerate}
\def\labelenumi{\arabic{enumi}.}
\tightlist
\item
  全体系の二乗ゼロ閉包を破らない。
\item
  均質背景から統計的・スペクトル的に識別できる。
\item
  安定した相対重心を持つ。
\item
  ゼロでも無限でもない有限広がりを持つ。
\item
  凝縮中心の外側に背景状態の変形を誘起し得る。
\item
  遠方または十分粗い読出しでは元の均質真空へ戻る。
\end{enumerate}

この六条件を満たす構造を、本稿では「局在凝縮体」と呼ぶ。

\begin{center}\rule{0.5\linewidth}{0.5pt}\end{center}

\subsection{2.
二乗ゼロ閉包のスケール不変性と干渉閉包}\label{ux4e8cux4e57ux30bcux30edux9589ux5305ux306eux30b9ux30b1ux30fcux30ebux4e0dux5909ux6027ux3068ux5e72ux6e09ux9589ux5305}

各複素数を

\[
z_n=a_n+i b_n
\]

と書くと、

\[
z_n^2
=
a_n^2-b_n^2
+
2i a_n b_n.
\]

したがって式 (1) は

\[
\sum_n(a_n^2-b_n^2)=0
\tag{2}
\]

および

\[
\sum_n a_n b_n=0
\tag{3}
\]

に分かれる。

特に式 (3)
は、各複素成分の交差項、すなわち二乗読出しに現れる干渉項の総和が閉じることを要求する。

全成分を共通実スケール \[\kappa>0\] で変換しても、

\[
a_n\mapsto \kappa a_n,
\qquad
b_n\mapsto \kappa b_n
\]

なら

\[
\sum_n(\kappa a_n)(\kappa b_n)
=
\kappa^2
\sum_n a_n b_n
=
0.
\]

したがってゼロ閉包はスケールを固定せず、

\[
\boxed{
\text{自己無撞着性は絶対スケールに依存しない}
}
\]

という性質を持つ。

本稿ではこの性質を重要視する。質量や長さを裸の絶対パラメータとして基礎方程式へ挿入するのではなく、凝縮後に生じる相対的な読出しとして扱う。

\begin{center}\rule{0.5\linewidth}{0.5pt}\end{center}

\subsection{3.
均質真空状態ゲージ場}\label{ux5747ux8ceaux771fux7a7aux72b6ux614bux30b2ux30fcux30b8ux5834}

\subsubsection{3.1 真空の定義}\label{ux771fux7a7aux306eux5b9aux7fa9}

真空状態を

\[
\mathcal V_0
=
\{z_n^{(0)}\}
\]

と書き、

\[
\sum_n
\left(z_n^{(0)}\right)^2
=
0
\tag{4}
\]

を満たすものとする。

さらに真空は、採用する匿名読出しに対して局在した特別点を持たないとする。この意味で「均質」と呼ぶ。

ここでいう状態ゲージとは、絶対的な共通振幅と共通位相を物理的に固定しないという意味である。すなわち

\[
z_n^{(0)}
\mapsto
\gamma z_n^{(0)}
\]

で同じ閉包クラスに属する状態は、相対読出しの観点では同一ゲージ軌道に属する。

\subsubsection{3.2
真空と物質を別の基礎変数にしない}\label{ux771fux7a7aux3068ux7269ux8ceaux3092ux5225ux306eux57faux790eux5909ux6570ux306bux3057ux306aux3044}

本稿では、真空と凝縮体を別種類の実体として導入しない。

凝縮後の状態を

\[
\mathcal V
=
\{z_n\}
\]

とし、

\[
\sum_n z_n^2=0
\tag{5}
\]

を維持する。

物質的凝縮とは、\[\mathcal V_0\] から \[\mathcal V\]
への自己無撞着な再配置のうち、局在した相関構造が生じた場合を指す。

したがって、

\[
\boxed{
\text{真空}
\longrightarrow
\text{同じ状態自由度の局在再配置}
\longrightarrow
\text{凝縮体}
}
\]

という一層構造を採用する。

\begin{center}\rule{0.5\linewidth}{0.5pt}\end{center}

\subsection{4.
無名集合の経験測度}\label{ux7121ux540dux96c6ux5408ux306eux7d4cux9a13ux6e2cux5ea6}

波の番号は物理的意味を持たないため、読出し演算は添字の並べ替えに依存してはならない。

有限集合の場合、経験測度を

\[
\mu
=
\frac1N
\sum_{n=1}^{N}
\delta_{z_n}
\tag{6}
\]

と定義する。

式 (6) は集合の並べ替えに不変である。

複素数を極形式

\[
z_n
=
r_n e^{i\theta_n},
\qquad
r_n>0
\]

で表す。

零点は対数振幅を持たないため、必要なら \[z=0\]
を極限点として別扱いし、変換は \[\mathbb C^\times\] 上で定義する。

非零複素数の乗法群は

\[
\mathbb C^\times
\simeq
\mathbb R_+
\times
S^1
\]

と分解できる。

したがって、振幅方向 \[r\in\mathbb R_+\] と位相方向 \[\theta\in S^1\]
に対する調和解析を組み合わせるのが自然である。

前報 {[}K1{]} の対数尺度読出し

\[
X=\ln r
\]

を採用すると、乗法的振幅比は加法的差へ写る。

\begin{center}\rule{0.5\linewidth}{0.5pt}\end{center}

\subsection{5. Anonymous Fourier--Mellin
変換}\label{anonymous-fouriermellin-ux5909ux63db}

\subsubsection{5.1 定義}\label{ux5b9aux7fa9}

無名複素集合に対して

\[
\boxed{
\mathcal A(s,m)
=
\frac1N
\sum_{n=1}^{N}
\exp
\left[
-i
\left(
s\ln r_n
+
m\theta_n
\right)
\right]
}
\tag{7}
\]

を定義する。

ここで

\[
s\in\mathbb R,
\qquad
m\in\mathbb Z
\]

である。

位相方向は円 \[S^1\] なので Fourier 級数、振幅方向は乗法群
\[\mathbb R_+\] なので Mellin 型変換になる。

式 (7) は経験測度を用いれば

\[
\mathcal A(s,m)
=
\int_{\mathbb C^\times}
|z|^{-is}
\left(
\frac{z}{|z|}
\right)^{-m}
d\mu(z)
\tag{8}
\]

と書ける。

\subsubsection{5.2 置換不変性}\label{ux7f6eux63dbux4e0dux5909ux6027}

任意の置換 \[\pi\] に対して

\[
z_n\mapsto z_{\pi(n)}
\]

としても総和は変わらない。したがって \[\mathcal A\] は無名性を保存する。

\subsubsection{5.3
共通スケール変換}\label{ux5171ux901aux30b9ux30b1ux30fcux30ebux5909ux63db}

\[
r_n
\mapsto
\kappa r_n
\]

なら、

\[
\ln r_n
\mapsto
\ln r_n+\ln\kappa.
\]

よって

\[
\mathcal A(s,m)
\mapsto
 e^{-is\ln\kappa}
\mathcal A(s,m).
\tag{9}
\]

したがって強度

\[
\boxed{
P(s,m)
=
|\mathcal A(s,m)|^2
}
\tag{10}
\]

は共通スケールに不変である。

\subsubsection{5.4
共通位相回転}\label{ux5171ux901aux4f4dux76f8ux56deux8ee2}

\[
\theta_n
\mapsto
\theta_n+\phi
\]

なら、

\[
\mathcal A(s,m)
\mapsto
 e^{-im\phi}
\mathcal A(s,m).
\tag{11}
\]

したがって式 (10) は共通位相にも不変である。

以上から \[P(s,m)\]
は、波の番号、絶対振幅スケール、絶対位相の三つに依存しない匿名スペクトルとなる。

Fourier--Mellin
型解析がスケール変換や回転に対する不変表現に利用できること自体は既知である
{[}E1{]}。本稿の独自点は、それを画像や既存空間座標に対して用いるのではなく、時空以前の無名複素状態集合に適用し、凝縮体の物質的読出しへ接続することである。

\begin{center}\rule{0.5\linewidth}{0.5pt}\end{center}

\subsection{6. 凝縮体の定義}\label{ux51ddux7e2eux4f53ux306eux5b9aux7fa9}

真空状態 \[\mu_0\] に対して凝縮後の測度を \[\mu\] とする。

形式的差を

\[
\delta\mu
=
\mu-\mu_0
\tag{12}
\]

と置く。

凝縮体とは、\[\delta\mu\]
の中に、背景から分離可能で、かつ有限幅を持つ相関成分が存在する状態である。

ここでは凝縮体成分を \[\mu_C\] と書き、その総重みを

\[
M_0
=
\int d\mu_C
\tag{13}
\]

とする。

\(M_0\)
はまだ物理質量とは呼ばず、「凝縮参加量」または「零次凝縮モーメント」と呼ぶ。

正規化した凝縮体測度を

\[
d\nu_C
=
\frac{d\mu_C}{M_0}
\tag{14}
\]

とする。

\begin{center}\rule{0.5\linewidth}{0.5pt}\end{center}

\subsection{7.
重心はスペクトル位相の一次微分として現れる}\label{ux91cdux5fc3ux306fux30b9ux30daux30afux30c8ux30ebux4f4dux76f8ux306eux4e00ux6b21ux5faeux5206ux3068ux3057ux3066ux73feux308cux308b}

\subsubsection{7.1
一次元対数尺度での導出}\label{ux4e00ux6b21ux5143ux5bfeux6570ux5c3aux5ea6ux3067ux306eux5c0eux51fa}

まず位相方向を固定し、対数尺度

\[
X=\ln r
\]

のみを見る。

凝縮体の正規化特性関数を

\[
\chi_C(s)
=
\int
 e^{-isX}
d\nu_C(X)
\tag{15}
\]

とする。

原点で

\[
\chi_C(0)=1.
\]

微分すると

\[
\left.
\frac{d\chi_C}{ds}
\right|_{s=0}
=
-i
\int X\,d\nu_C
=
-iX_C,
\tag{16}
\]

ここで

\[
\boxed{
X_C
=
\int X\,d\nu_C
}
\tag{17}
\]

を凝縮体の重心と定義する。

特性関数の対数を用いると、

\[
\boxed{
X_C
=
i
\left.
\frac{d}{ds}
\ln\chi_C(s)
\right|_{s=0}
}
\tag{18}
\]

となる。

符号規約を位相で書けば、

\[
\boxed{
X_C
=
-
\left.
\frac{d}{ds}
\arg\chi_C(s)
\right|_{s=0}
}
\tag{19}
\]

である。

これは、凝縮体の中心位置がスペクトル位相の一次傾きとして読めることを示す。

\subsubsection{7.2
絶対重心は読出さない}\label{ux7d76ux5bfeux91cdux5fc3ux306fux8aadux51faux3055ux306aux3044}

全凝縮体を

\[
X\mapsto X+X_0
\]

と共通移動すると、

\[
\chi_C(s)
\mapsto
 e^{-isX_0}\chi_C(s).
\tag{20}
\]

したがって \[|\chi_C|\] は変わらないが、位相傾きは \[X_0\]
だけ変化する。

これは欠点ではない。基礎系が絶対スケールを持たない以上、絶対重心も物理的に固定されるべきではない。

二つの凝縮体 \[A,B\] に対して、

\[
\frac{\chi_A(s)}{\chi_B(s)}
\]

を取ると、

\[
\boxed{
X_A-X_B
=
-
\left.
\frac{d}{ds}
\arg
\frac{\chi_A(s)}{\chi_B(s)}
\right|_{s=0}
}
\tag{21}
\]

が得られる。

したがって読出し可能なのは絶対重心ではなく相対重心である。

\begin{center}\rule{0.5\linewidth}{0.5pt}\end{center}

\subsection{8.
有限広がりはスペクトル振幅の二次微分として現れる}\label{ux6709ux9650ux5e83ux304cux308aux306fux30b9ux30daux30afux30c8ux30ebux632fux5e45ux306eux4e8cux6b21ux5faeux5206ux3068ux3057ux3066ux73feux308cux308b}

凝縮体の分散を

\[
\sigma_X^2
=
\int
(X-X_C)^2
d\nu_C
\tag{22}
\]

とする。

特性関数のキュムラント展開 {[}E2{]} により、

\[
\ln\chi_C(s)
=
-isX_C
-
\frac12
s^2\sigma_X^2
+
O(s^3).
\tag{23}
\]

したがって

\[
\boxed{
\sigma_X^2
=
-
\left.
\frac{d^2}{ds^2}
\Re\ln\chi_C(s)
\right|_{s=0}
}
\tag{24}
\]

となる。

これは、

\[
\boxed{
\text{一次スペクトル位相}
\longrightarrow
\text{重心}
}
\]

\[
\boxed{
\text{二次スペクトル減衰}
\longrightarrow
\text{広がり}
}
\]

という明瞭な階層を与える。

物質的凝縮体の最低条件として

\[
0<\sigma_X^2<\infty
\tag{25}
\]

を要求すれば、完全な点特異点でも無限に拡散した背景でもない有限サイズの対象を定義できる。

\begin{center}\rule{0.5\linewidth}{0.5pt}\end{center}

\subsection{9.
位相・時間方向を含む二変数展開}\label{ux4f4dux76f8ux6642ux9593ux65b9ux5411ux3092ux542bux3080ux4e8cux5909ux6570ux5c55ux958b}

位相方向も含め、

\[
Y=
\begin{pmatrix}
X\\
T
\end{pmatrix}
\]

とする。

\(T\) は前報 {[}K1{]}
における位相周期または相対巻き数から得られる時間読出しである。

双対変数を

\[
q=
\begin{pmatrix}
s\\
\omega
\end{pmatrix}
\]

とすると、

\[
\chi_C(q)
=
\int
 e^{-iq^{\mathsf T}Y}
d\nu_C(Y).
\tag{26}
\]

原点近傍で

\[
\boxed{
\ln\chi_C(q)
=
-iq^{\mathsf T}\mu_C
-
\frac12
q^{\mathsf T}
\Sigma_C
q
+
O(|q|^3)
}
\tag{27}
\]

となる。

ここで

\[
\mu_C
=
\begin{pmatrix}
X_C\\
T_C
\end{pmatrix}
\]

は重心、

\[
\Sigma_C
=
\begin{pmatrix}
\sigma_X^2 & C_{XT}\\
C_{XT} & \sigma_T^2
\end{pmatrix}
\]

は共分散行列である。

したがって交差項

\[
\boxed{
C_{XT}
=
-
\left.
\frac{\partial^2}
{\partial s\,\partial\omega}
\Re\ln\chi_C
\right|_{q=0}
}
\tag{28}
\]

も読出せる。

これにより、凝縮体の単なる半径だけでなく、空間尺度と時間位相の結合も測れる。

\begin{center}\rule{0.5\linewidth}{0.5pt}\end{center}

\subsection{10.
3空間＋時間読出しへの一般化}\label{ux7a7aux9593ux6642ux9593ux8aadux51faux3057ux3078ux306eux4e00ux822cux5316}

\textbf{仮定 10.1（4次元相対読出しの存在）}\\
前報 {[}K1{]}
の時空読出し、および複数複素平面の相対幾何を拡張し、凝縮体に対して4次元の相対読出し座標が構成できると仮定する。これは本稿で新たに証明する命題ではない。前報で明示的に構成した読出しはまず
\[1+1\]
次元であり、以下の4次元化はその拡張が成立する場合の条件付き構成である。

具体的に、凝縮体に対して読出し座標

\[
Y_n
=
\begin{pmatrix}
X_n\\
Y_n\\
Z_n\\
T_n
\end{pmatrix}
\]

が構成できたとする。

ここで重要なのは、これらが最初から背景格子として与えられているのではなく、複素状態関係から得られた相対読出しであることである。

従って、式 (29)〜(31) の4次元モーメント読出し、およびそれを用いる §11
の4次元的な質量的解釈は仮定 10.1 に条件付きである。一方、§7〜§9
で導出した一次元・二変数のモーメント恒等式そのものは、この仮定に依存しない。

4次元双対変数を

\[
q
=
\begin{pmatrix}
k_x\\
k_y\\
k_z\\
\omega
\end{pmatrix}
\]

とし、

\[
\boxed{
\chi_C(q)
=
\frac{1}{M_0}
\int
\exp
\left[
-i
\left(
k_xX+k_yY+k_zZ-\omega T
\right)
\right]
d\mu_C
}
\tag{29}
\]

とする。

符号は Lorentz
型読出しとの整合のため時間項に反対符号を採用したが、モーメント抽出自体は符号規約に依存しない。

式 (29) の対数展開は

\[
\boxed{
\ln\chi_C(q)
=
-iq^{\mathsf T}\mu_C
-
\frac12q^{\mathsf T}\Sigma_C q
+
O(|q|^3)
}
\tag{30}
\]

である。

ここで

\[
\boxed{
\mu_C
=
(X_C,Y_C,Z_C,T_C)
}
\tag{31}
\]

が4次元相対重心であり、\[\Sigma_C\] は \[4\times4\] の共分散行列である。

空間部分 \[3\times3\]
の固有値・固有ベクトルから、主半径、楕円率、主軸方向、異方性を読出せる。

したがって単一のスペクトル演算から、凝縮量、重心、広がり、形状、時空間相関を統一的に得られる。

\begin{center}\rule{0.5\linewidth}{0.5pt}\end{center}

\subsection{11. 質量的計量}\label{ux8ceaux91cfux7684ux8a08ux91cf}

\subsubsection{11.1
質量を基礎パラメータとして置かない}\label{ux8ceaux91cfux3092ux57faux790eux30d1ux30e9ux30e1ux30fcux30bfux3068ux3057ux3066ux7f6eux304bux306aux3044}

本稿では物理質量 \[m\] を初期条件として導入しない。

均質真空には局在中心がない。

凝縮が起こると初めて、凝縮参加量、相対重心、有限広がり、背景歪みが同時に定義可能になる。

そこで質量的対象の最小データを

\[
\boxed{
\mathfrak M_C
=
\left(
M_0,
\mu_C,
\Sigma_C
\right)
}
\tag{32}
\]

と定義する。

これはまだ kg 単位の慣性質量ではない。

本稿でいう「質量的計量」とは、均質真空に対して局在した中心と有限幅を作る凝縮強度を表す状態量である。

\subsubsection{11.2
零次モーメントとしての凝縮量}\label{ux96f6ux6b21ux30e2ux30fcux30e1ux30f3ux30c8ux3068ux3057ux3066ux306eux51ddux7e2eux91cf}

非正規化凝縮スペクトルを

\[
\Delta\mathcal A_C(q)
=
\int
 e^{-iq^{\mathsf T}Y}
d\mu_C(Y)
\tag{33}
\]

とすると、

\[
\boxed{
\Delta\mathcal A_C(0)
=
M_0
}
\tag{34}
\]

である。

したがって、

\[
\boxed{
\begin{aligned}
0\text{次モーメント}
&\longrightarrow
\text{凝縮量},\\
1\text{次モーメント}
&\longrightarrow
\text{重心},\\
2\text{次モーメント}
&\longrightarrow
\text{広がり・形状}
\end{aligned}
}
\tag{35}
\]

という階層が成立する。

この構造は、質量を外部スカラーとして付与するのではなく、凝縮体の状態分布そのものから計量するための最小候補を与える。

\begin{center}\rule{0.5\linewidth}{0.5pt}\end{center}

\subsection{12.
均質真空に対するクロス位相と相対位置}\label{ux5747ux8ceaux771fux7a7aux306bux5bfeux3059ux308bux30afux30edux30b9ux4f4dux76f8ux3068ux76f8ux5bfeux4f4dux7f6e}

スペクトル強度

\[
|\mathcal A|^2
\]

だけでは絶対位置情報は失われる。

したがって、凝縮体の重心を読むには参照とのクロス位相が必要である。

凝縮前真空を \[\mathcal A_0(q)\]、凝縮後を \[\mathcal A(q)\] とする。

差スペクトル

\[
\boxed{
\delta\mathcal A(q)
=
\mathcal A(q)-\mathcal A_0(q)
}
\tag{36}
\]

を考える。

あるいは、参照スペクトルが非零の領域ではクロススペクトル

\[
\boxed{
C(q)
=
\mathcal A(q)
\mathcal A_0(q)^*
}
\tag{37}
\]

を用いてもよい。

ここで得られる位相は、外部原点に対する位置ではなく、均質真空ゲージに対する相対位相ずれである。

この構造は、本稿の基礎原則「絶対量ではなく関係だけを読む」と整合する。

\begin{center}\rule{0.5\linewidth}{0.5pt}\end{center}

\subsection{13.
凝縮体による周囲真空状態ゲージの歪み}\label{ux51ddux7e2eux4f53ux306bux3088ux308bux5468ux56f2ux771fux7a7aux72b6ux614bux30b2ux30fcux30b8ux306eux6b6aux307f}

\subsubsection{13.1
局在凝縮と全体閉包}\label{ux5c40ux5728ux51ddux7e2eux3068ux5168ux4f53ux9589ux5305}

全体系では常に

\[
\sum_n z_n^2=0
\]

を維持する。

凝縮体領域 \[C\] とその外部 \[E\] を形式的に分けると、

\[
Q_C
=
\sum_{n\in C}z_n^2,
\]

\[
Q_E
=
\sum_{n\in E}z_n^2
\]

であり、

\[
\boxed{
Q_C+Q_E=0
}
\tag{38}
\]

である。

一般に凝縮体部分だけが \[Q_C=0\] である必要はない。

もし局在凝縮が

\[
Q_C\neq0
\]

を持つなら、全体閉包は必ず

\[
\boxed{
Q_E=-Q_C
}
\tag{39}
\]

を要求する。

したがって、局所凝縮によって閉包寄与が偏れば、周囲の状態集合は補償的に再配置されなければならない。

これは、

\[
\boxed{
\text{局所凝縮}
\longrightarrow
\text{背景状態の補償変形}
}
\]

を追加公理なしで生じさせる候補である。

\subsubsection{13.2
可約閉包と残差補償機構の限界}\label{ux53efux7d04ux9589ux5305ux3068ux6b8bux5deeux88dcux511fux6a5fux69cbux306eux9650ux754c}

前報 {[}K1{]}
の明示的殻特殊解は、各殻がそれぞれ単独で二乗ゼロ閉包する可約構造を持つ。このため、完全な閉包殻だけを集めて凝縮体領域
\[C\] を構成した場合には

\[
Q_C=0
\]

となり、式 (39) の \textbf{非零閉包残差を外部が補償する機構}
は作動しない。従って、実験によって式 (39)
型の補償機構そのものを検証するには、凝縮体が閉包単位を横断して切り出され、一般に
\[Q_C\neq0\] となり得る構成を含める必要がある。

ただし、これは

\[
Q_C=0
\]

なら周囲の背景変形が一切存在しないことを意味しない。二乗和の総閉包を保ったまま、内部相関、位相配置、あるいは凝縮体と外部との関係構造が再配置される可能性は残る。従って本稿では、

\begin{itemize}
\tightlist
\item
  \textbf{閉包残差補償型の歪み}：\[Q_C\neq0\] に対して式 (39)
  が要求する変形、
\item
  \textbf{閉包保存型の相関歪み}：\[Q_C=0\]
  を保ちながら生じ得る別種の再配置、
\end{itemize}

を区別する。後者の存在条件と読出しは未解決問題である。

\subsubsection{13.3
歪みのスペクトル読出し}\label{ux6b6aux307fux306eux30b9ux30daux30afux30c8ux30ebux8aadux51faux3057}

凝縮前後の差 \[\delta\mathcal A(q)\] を逆変換して \[\delta\rho(Y)\]
を得る。

概念的に、

\[
\delta\rho(Y)
=
\mathcal F^{-1}
\left[
\delta\mathcal A(q)
\right]
\tag{40}
\]

である。

ただし空間尺度が対数振幅から得られる場合、元の複素状態変数に対しては
Fourier 逆変換と Mellin 逆変換の積として解釈される。

局在凝縮体が真に周囲のゲージを歪めるなら、\[\delta\rho(Y)\neq0\]
が凝縮中心の外側にも存在し、

\[
\delta\rho(Y)
\to0
\qquad
\text{for sufficiently large relative separation}
\tag{41}
\]

となることが期待される。

このとき、凝縮中心、有限コア、外部へ減衰する背景変形という孤立物体に必要な構造が得られる。

\begin{center}\rule{0.5\linewidth}{0.5pt}\end{center}

\subsection{14.
背景曲率との関係}\label{ux80ccux666fux66f2ux7387ux3068ux306eux95a2ux4fc2}

前報 {[}K1{]}
では、閉じた有界な複素状態場から加法的な読出し空間を構成する可能性を扱った。

本稿では、凝縮体による状態分布変形を \[\delta\rho\] として抽出する。

読出し計量を \[g_0\]、凝縮後を \[g_C\] と書くなら、

\[
\boxed{
\delta g
=
g_C-g_0
}
\tag{42}
\]

を凝縮による局所的計量偏差と定義することが将来的に可能である。

ただし、本稿では \[\delta\rho\] から \[\delta g\]
への一意な写像はまだ導出しない。

したがって本稿で確立するのは、

\[
\boxed{
\text{局在凝縮}
\longrightarrow
\text{真空状態分布の非均質}
}
\]

までであり、

\[
\boxed{
\delta\rho
\longrightarrow
\text{Einstein 曲率}
}
\]

は未解決問題として残す。

この区別は重要である。

一方、凝縮による周囲変形が存在するかどうかは、数値実験によって直接検証可能である。

\begin{center}\rule{0.5\linewidth}{0.5pt}\end{center}

\subsection{15. 4次元 FFT
との関係}\label{ux6b21ux5143-fft-ux3068ux306eux95a2ux4fc2}

見かけ上、式 (29) は4次元 Fourier 変換に近い。

しかし通常の4次元 FFT では、\[ (x,y,z,t) \] の格子が最初から存在する。

本稿では順序が逆である。

\[
\boxed{
\text{無名複素波}
\rightarrow
\text{相対尺度・位相読出し}
\rightarrow
\text{創発座標}
\rightarrow
\text{スペクトル凝縮解析}
}
\tag{43}
\]

である。

特に、

\[
X=\ln r
\]

であるため、

\[
e^{-isX}
=
e^{-is\ln r}
=
r^{-is}
\]

となる。

したがって元の振幅変数 \[r\] に対しては Mellin 核そのものである。

よって本稿の演算は単なる4次元 FFT ではなく、

\[
\boxed{
\text{創発対数空間における Fourier}
\equiv
\text{元振幅空間における Mellin}
}
\]

という対応を持つ。

\begin{center}\rule{0.5\linewidth}{0.5pt}\end{center}

\subsection{16.
無名状態からの基底探索}\label{ux7121ux540dux72b6ux614bux304bux3089ux306eux57faux5e95ux63a2ux7d22}

ここまでの式は、凝縮体に対応する読出し座標が得られた後のモーメント抽出を与える。

しかし最も難しい問題は、その基底自体を無名状態からどう選ぶかである。

本稿では、候補基底 \[B\] に対するスペクトル

\[
\mathcal A_B(q)
\]

を考え、凝縮度評価関数

\[
\mathcal C(B)
\]

を定義することを提案する。

例えば、

\[
\mathcal C(B)
=
\frac{
\displaystyle
\int_{\Omega_{\rm peak}}
|\delta\mathcal A_B(q)|^2dq
}{
\displaystyle
\int_{\Omega_{\rm all}}
|\delta\mathcal A_B(q)|^2dq
}
\tag{44}
\]

のように、限られたスペクトル領域への集中度を測ることができる。

最適基底を

\[
\boxed{
B_*
=
\operatorname*{arg\,max}_{B}
\mathcal C(B)
}
\tag{45}
\]

とする。

ただし許される基底 \[B\]
は任意ではなく、置換不変な状態量だけから構成されること、共通スケール・共通位相ゲージに依存しないこと、二乗ゼロ閉包を保存すること、前報
{[}K1{]} の時空読出しと整合することを要求する。

したがって、これは座標を自由にフィットして任意の凝縮を作る操作ではない。

むしろ、

\[
\boxed{
\text{ゼロ閉包系自身が許す基底の中で、
最も自己無撞着に局在するモードを探す}
}
\]

問題である。

\begin{center}\rule{0.5\linewidth}{0.5pt}\end{center}

\subsection{17.
数値実験としての最小検証}\label{ux6570ux5024ux5b9fux9a13ux3068ux3057ux3066ux306eux6700ux5c0fux691cux8a3c}

本稿の提案は数値的に直接検証できる。

\subsubsection{実験 A：匿名性}\label{ux5b9fux9a13-aux533fux540dux6027}

既知の凝縮体を含む複素波集合を作り、全波の番号をランダムにシャッフルする。

判定条件は

\[
\mathcal A_{\rm before}
=
\mathcal A_{\rm after}
\]

である。

\subsubsection{実験
B：共通スケール不変性}\label{ux5b9fux9a13-bux5171ux901aux30b9ux30b1ux30fcux30ebux4e0dux5909ux6027}

全波を

\[
z_n\mapsto\kappa z_n
\]

と変換する。

判定条件は

\[
|\mathcal A_{\kappa}(s,m)|^2
=
|\mathcal A(s,m)|^2.
\]

\subsubsection{実験
C：共通位相不変性}\label{ux5b9fux9a13-cux5171ux901aux4f4dux76f8ux4e0dux5909ux6027}

\[
z_n\mapsto e^{i\phi}z_n
\]

と変換する。

判定条件は同様に

\[
|\mathcal A_{\phi}(s,m)|^2
=
|\mathcal A(s,m)|^2.
\]

\subsubsection{実験
D：重心回収}\label{ux5b9fux9a13-dux91cdux5fc3ux56deux53ce}

既知の相対移動 \[\Delta X\] を凝縮体へ加える。

予測は

\[
\chi_C(s)
\mapsto
 e^{-is\Delta X}
\chi_C(s)
\]

であり、

\[
{-}
\left.
\frac{d}{ds}
\Delta\arg\chi_C
\right|_0
=
\Delta X
\tag{46}
\]

が回収されるはずである。

\subsubsection{実験
E：有限幅回収}\label{ux5b9fux9a13-eux6709ux9650ux5e45ux56deux53ce}

既知の分散 \[\sigma^2\] を持つ凝縮体を生成する。

式 (24) により回収した分散と入力値を比較する。

\subsubsection{実験
F：複数凝縮体}\label{ux5b9fux9a13-fux8907ux6570ux51ddux7e2eux4f53}

二つ以上の凝縮体を配置し、\[X_A-X_B\] をクロス位相から回収する。

絶対原点を変えても結果が不変であることを確認する。

\subsubsection{実験
G：背景歪み}\label{ux5b9fux9a13-gux80ccux666fux6b6aux307f}

全体ゼロ閉包を保ったまま局所凝縮を発生させ、\[\delta\rho(Y)\]
が凝縮コアの外側へどの程度広がるかを測定する。特に式 (39)
型の残差補償機構を検証する系列では、完全な閉包殻だけを凝縮体として選ばず、閉包単位を横断して一般に
\[Q_C\neq0\] となり得る凝縮を生成条件に含める。これと並行して、\[Q_C=0\]
を維持する凝縮でも閉包保存型の相関歪みが生じるかを別系列として測る。

これが本稿で最も物理的に重要な検証である。

\subsubsection{実験
H：ブラインド対照と偽陽性率}\label{ux5b9fux9a13-hux30d6ux30e9ux30a4ux30f3ux30c9ux5bfeux7167ux3068ux507dux967dux6027ux7387}

凝縮体を一切植え込まない均質真空状態を対照群として生成し、§16
の基底探索を同じ条件で実行する。候補基底ごとの集中度
\[\mathcal C(B)\]、差スペクトル
\[\delta\mathcal A\]、見かけの重心・有限幅、および背景歪み指標の分布を測定し、凝縮を含まない系でも生じる偽陽性率を先に確定する。

その上で実験 D〜G
の信号が、この対照分布から統計的に分離できることを要求する。これにより、基底探索が凝縮を発見しているのか、それとも最適化によって見かけの局在を作っているのかを判別する。

\begin{center}\rule{0.5\linewidth}{0.5pt}\end{center}

\subsection{18.
繰り込みとの関係}\label{ux7e70ux308aux8fbcux307fux3068ux306eux95a2ux4fc2}

本稿の基礎方程式

\[
\sum_n z_n^2=0
\]

は共通スケール変換に対して不変である。

したがって、基礎状態の自己無撞着性を保つために特定の絶対スケールを選ぶ必要はない。

また質量的計量を裸の質量パラメータとして投入せず、凝縮量、相対重心、有限幅、背景変形として読出す。

このため、本構成は少なくとも「発散する裸パラメータを観測量へ再調整すること」を基礎定義として必要としない。

ただし、この事実だけから量子場理論の全ての紫外発散が自動的に消えるとはまだ結論しない。

本稿の厳密な主張は、

\[
\boxed{
\text{基礎閉包条件はスケール不変であり、
質量的量を外部の裸パラメータとして導入しない}
}
\]

という点までである。

将来、凝縮体間相互作用を導入した際にも全観測量が有限に保たれることが示されれば、そこで初めてより強い
UV-finite 性を議論できる。

\begin{center}\rule{0.5\linewidth}{0.5pt}\end{center}

\subsection{19.
標準的な質量概念との違い}\label{ux6a19ux6e96ux7684ux306aux8ceaux91cfux6982ux5ff5ux3068ux306eux9055ux3044}

本稿の「質量的計量」は、標準模型の Higgs
機構、慣性質量、重力質量を既に同一視するものではない。

本稿で構成したのは、

\[
\boxed{
\text{均質真空の中に固有重心を生じさせる局在凝縮の尺度}
}
\]

である。

したがって今後区別して検証すべき量は少なくとも、凝縮参加量
\[M_0\]、慣性応答としての質量、周囲背景歪みの強度、重力的遠方場の係数、内部振動周期または固有周波数である。

もし最終的にこれらが同一の凝縮不変量へ収束すれば、質量の統一的読出しとなる。

現段階ではその一致を仮定しない。

\begin{center}\rule{0.5\linewidth}{0.5pt}\end{center}

\subsection{20.
本構成の最小性}\label{ux672cux69cbux6210ux306eux6700ux5c0fux6027}

本稿で追加した数学的構造は限定的である。

基礎として必要なのは、

\[
\boxed{
\begin{aligned}
&\text{(i) 無名な複素状態集合},\\
&\text{(ii) 二乗ゼロ閉包},\\
&\text{(iii) 前報の相対対数尺度・位相読出し},\\
&\text{(iv) 経験測度に対する調和解析}
\end{aligned}
}
\]

である。

粒子ラベル、外部座標格子、絶対原点、裸質量は導入しない。

また、局在体の中心を先に指定して Fourier
変換するのではない。中心は変換後の一次モーメントから出る。

したがって論理順序は

\[
\boxed{
\text{波}
\rightarrow
\text{凝縮}
\rightarrow
\text{スペクトル}
\rightarrow
\text{重心}
}
\]

であり、既知の位置に波束を置くという順序ではない。

\begin{center}\rule{0.5\linewidth}{0.5pt}\end{center}

\subsection{21. 未解決問題}\label{ux672aux89e3ux6c7aux554fux984c}

\subsubsection{21.1
凝縮体は自発的に生成されるか}\label{ux51ddux7e2eux4f53ux306fux81eaux767aux7684ux306bux751fux6210ux3055ux308cux308bux304b}

本稿は凝縮体を読出す方法を与えたが、

\[
\sum_n z_n^2=0
\]

だけから安定した局在凝縮が必ず生成されることは証明していない。

これは最重要の動力学的問題である。

\subsubsection{21.2
凝縮体の安定性}\label{ux51ddux7e2eux4f53ux306eux5b89ux5b9aux6027}

重心 \[\mu_C\] と共分散 \[\Sigma_C\]
が長時間安定する条件を求める必要がある。

\subsubsection{21.3
凝縮量と背景歪みの関係}\label{ux51ddux7e2eux91cfux3068ux80ccux666fux6b6aux307fux306eux95a2ux4fc2}

\[
M_0
\]

と

\[
\delta\rho
\]

の遠方振幅の間に単純な比例則があるかを検証する必要がある。

\subsubsection{21.4 逆二乗則}\label{ux9006ux4e8cux4e57ux5247}

3次元読出し後、孤立凝縮体の背景歪みが遠方で

\[
\frac{1}{R^2}
\]

型の勾配を生むかは未検証である。

\subsubsection{21.5
慣性と重力の一致}\label{ux6163ux6027ux3068ux91cdux529bux306eux4e00ux81f4}

凝縮体の加速抵抗と周囲背景歪みが同一不変量で決まるなら、慣性質量と重力質量の一致に接続できる可能性がある。現段階では予測課題である。

\subsubsection{21.6
凝縮体分解の一意性}\label{ux51ddux7e2eux4f53ux5206ux89e3ux306eux4e00ux610fux6027}

本稿では差測度 \[\delta\mu\] から凝縮体成分 \[\mu_C\]
を分離して重心と広がりを読む。しかし、無名な多体系に複数の局在相関が共存するとき、

\[
\delta\mu
\longrightarrow
\{\mu_{C_1},\mu_{C_2},\ldots\}
\]

という分解が一意に定まるとはまだ示していない。§16
の基底探索における最大集中度だけでなく、分解の存在、一意性、近接凝縮体の分離限界、および分解能依存性を定式化する必要がある。これは「どこまでを一個の物体と呼ぶか」を外部から指定しないための中心課題である。

\subsubsection{21.7
閉包残差補償型と閉包保存型の背景歪み}\label{ux9589ux5305ux6b8bux5deeux88dcux511fux578bux3068ux9589ux5305ux4fddux5b58ux578bux306eux80ccux666fux6b6aux307f}

§13.2
で区別した二種類の背景変形が、どの条件で生じ、同一の遠方読出しへ収束するかは未解決である。特に、\[Q_C=0\]
の凝縮でも相関再配置だけから長距離歪みが生じ得るかを判別する必要がある。

\begin{center}\rule{0.5\linewidth}{0.5pt}\end{center}

\subsection{22.
議論――真空と物質の境界を別の実体にしない}\label{ux8b70ux8ad6ux771fux7a7aux3068ux7269ux8ceaux306eux5883ux754cux3092ux5225ux306eux5b9fux4f53ux306bux3057ux306aux3044}

本構成の特徴は、真空と物質を基礎的に別物として置かない点にある。

均質真空は

\[
\mathcal V_0
\]

であり、凝縮体を含む宇宙は

\[
\mathcal V
\]

である。

両者は同じ複素状態自由度からなり、違いは相関構造にある。

したがって、

\[
\boxed{
\text{物質}
=
\text{真空状態ゲージ場の局在凝縮}
}
\tag{47}
\]

という読みが可能になる。

さらに、凝縮によって全体ゼロ閉包を維持するための補償変形が周囲へ及ぶなら、物質と背景歪みは独立の入力ではなく、同じ自己無撞着な再配置の二つの読出しとなる。

ここに本構成の最も重要な物理的可能性がある。

\begin{center}\rule{0.5\linewidth}{0.5pt}\end{center}

\subsection{23. 結論}\label{ux7d50ux8ad6}

本稿では、無名な複素定常波からなる均質真空状態場

\[
\sum_n z_n^2=0
\]

の内部に局在凝縮体が存在する場合、その凝縮体の物質的情報を外部座標や既知質量なしに読出す枠組みを構成した。

主要結果は以下である。

第一に、無名複素集合の経験測度に対して

\[
\mathcal A(s,m)
=
\frac1N
\sum_n
\exp
\left[
-i
\left(
s\ln r_n+m\theta_n
\right)
\right]
\]

という Fourier--Mellin 型変換を定義した。

第二に、そのスペクトル強度は波の置換、共通振幅スケール、共通位相回転に不変である。

第三に、凝縮体の正規化変換の対数展開

\[
\ln\chi_C(q)
=
-iq^{\mathsf T}\mu_C
-
\frac12q^{\mathsf T}\Sigma_Cq
+
\cdots
\]

から、一次項として相対重心、二次項として有限広がり・形状を得た。

第四に、非正規化零次項を凝縮参加量とすることで、

\[
\boxed{
0\text{次}
\to
\text{凝縮量},
\qquad
1\text{次}
\to
\text{重心},
\qquad
2\text{次}
\to
\text{広がり}
}
\]

という質量的モーメント階層を構成した。

第五に、凝縮前後の差スペクトル

\[
\delta\mathcal A
=
\mathcal A-\mathcal A_0
\]

から、凝縮体の外部へ広がる真空状態ゲージの歪みを直接検出する方法を示した。

したがって、本構成では

\[
\boxed{
\text{均質真空}
\rightarrow
\text{局在凝縮}
\rightarrow
\text{相対重心}
\rightarrow
\text{有限広がり}
\rightarrow
\text{質量的計量}
\rightarrow
\text{周囲ゲージ歪み}
}
\]

を、同一の無名複素状態場の内部読出しとして連続的に記述できる。

最も重要な未解決問題は、読み出し方法ではなく、二乗ゼロ閉包した均質定常波真空の動力学から、このような有限幅の安定凝縮体が自発生成することを示せるかである。

その意味で本稿は「質量の完成理論」ではなく、

\[
\boxed{
\text{無名な真空波の海から、
物体を物体として読出すための最小数学的装置}
}
\]

を与えるものである。

\begin{center}\rule{0.5\linewidth}{0.5pt}\end{center}

\section{補遺A　自己双対真空と非 null
読出しへの最小拡張}\label{ux88dcux907aa-ux81eaux5df1ux53ccux5bfeux771fux7a7aux3068ux975e-null-ux8aadux51faux3057ux3078ux306eux6700ux5c0fux62e1ux5f35}

\subsection{A.1　目的と位置づけ}\label{a.1-ux76eeux7684ux3068ux4f4dux7f6eux3065ux3051}

前報 {[}K1{]}
の可算無限特殊解では、空間尺度と周期尺度が同じ殻パラメータに比例し、任意の殻対について写像後の波長--振動数積が一定となる。さらに読出し単位を一致させると、全ての殻対区間は
null となる。

本補遺では、この構造を壊さずに最小限一般化し、どの条件で
spacelike、null、timelike
の三分類が開くかを調べる。ここで導出するのは読出し尺度間の数学的分類であり、局在凝縮体が実際に
timelike 構造を生成することは最後に独立した仮説として区別する。

正の空間尺度と周期尺度をそれぞれ

\[
L_k>0,
\qquad
P_k>0
\]

とし、対数読出しを

\[
X_k=C_X\ln L_k,
\qquad
T_k=C_T\ln P_k
\tag{A1}
\]

とする。殻対 \[j,k\] に対して

\[
\Delta X_{jk}
=
C_X\ln\frac{L_j}{L_k},
\qquad
\Delta T_{jk}
=
C_T\ln\frac{P_j}{P_k}
\tag{A2}
\]

である。

波長読出しと周期・振動数読出しを

\[
\lambda_{jk}=|\Delta X_{jk}|,
\qquad
\tau_{jk}=|\Delta T_{jk}|,
\qquad
\nu_{jk}=\frac{1}{\tau_{jk}}
\tag{A3}
\]

と定義する。

\begin{center}\rule{0.5\linewidth}{0.5pt}\end{center}

\subsection{\texorpdfstring{A.2　全殻対で \(\lambda\nu\)
が一定となるための必要十分条件}{A.2　全殻対で \textbackslash lambda\textbackslash nu が一定となるための必要十分条件}}\label{a.2-ux5168ux6bbbux5bfeux3067-lambdanu-ux304cux4e00ux5b9aux3068ux306aux308bux305fux3081ux306eux5fc5ux8981ux5341ux5206ux6761ux4ef6}

\subsubsection{定理 A1}\label{ux5b9aux7406-a1}

少なくとも三つの相異なる空間尺度 \[L_k\]
を含む殻族を考える。全ての異なる殻対 \[j\neq k\] について

\[
\lambda_{jk}\nu_{jk}=K
\qquad
(K>0)
\tag{A4}
\]

が同じ定数となるための必要十分条件は、ある定数 \[P_0>0\] と非零実数
\[\alpha\] が存在して

\[
\boxed{
P_k=P_0L_k^{\alpha}
}
\tag{A5}
\]

が全殻で成立することである。軸向きの反転は \[\alpha\] の符号または
\[C_T\] の符号に吸収できるため、尺度の増加関係だけを扱う場合は
\[\alpha>0\] としてよい。

\subsubsection{証明}\label{ux8a3cux660e}

\[
x_k=\ln L_k,
\qquad
y_k=\ln P_k
\]

と置く。式 (A3)(A4) から

\[
\frac{|C_X|\,|x_j-x_k|}
{|C_T|\,|y_j-y_k|}
=K
\]

なので、ある \[a>0\] に対して

\[
|y_j-y_k|
=
a|x_j-x_k|
\tag{A6}
\]

が全殻対で成立する。

したがって写像 \[x_k\mapsto y_k/a\]
は、実数直線上のこの点集合の全ての対距離を保存する。三点以上を含む実数直線上の距離保存写像は、平行移動と向き反転を除いて一意であるため、ある
\[b\in\mathbb R\] と \[\sigma=\pm1\] により

\[
y_k
=
\sigma a x_k+b
\tag{A7}
\]

と書ける。ここで

\[
\alpha=\sigma a,
\qquad
P_0=e^b
\]

とすれば式 (A5) を得る。

逆に式 (A5) が成立すれば

\[
\ln\frac{P_j}{P_k}
=
\alpha
\ln\frac{L_j}{L_k}
\]

なので、式 (A2)(A3) より

\[
\boxed{
\lambda_{jk}\nu_{jk}
=
\frac{|C_X|}{|C_T|\,|\alpha|}
}
\tag{A8}
\]

となり、殻対に依存しない。以上で必要十分性が示された。□

この定理により、前報 {[}K1{]} の条件 \[P_k\propto L_k\]
は一般条件の特殊点

\[
\boxed{\alpha=1}
\]

であることが分かる。

\begin{center}\rule{0.5\linewidth}{0.5pt}\end{center}

\subsection{\texorpdfstring{A.3　\(\alpha\) による spacelike / null /
timelike
分類}{A.3　\textbackslash alpha による spacelike / null / timelike 分類}}\label{a.3-alpha-ux306bux3088ux308b-spacelike-null-timelike-ux5206ux985e}

読出し単位を

\[
|C_X|=|C_T|=:C
\tag{A9}
\]

に規格化する。式 (A5) から

\[
\Delta T_{jk}
=
\sigma\alpha\,\Delta X_{jk}
\]

となる。ここで \[\sigma=\pm1\]
は軸向きだけを表すので、二次形式には影響しない。

Lorentz 型読出し二次形式を

\[
ds_{jk}^2
=
(\Delta X_{jk})^2
-
(\Delta T_{jk})^2
\tag{A10}
\]

とすれば、

\[
\boxed{
 ds_{jk}^2
=
(1-\alpha^2)
C^2
\left(
\ln\frac{L_j}{L_k}
\right)^2
}
\tag{A11}
\]

を得る。

したがって \[\alpha>0\] の場合、

\[
\boxed{
0<\alpha<1
\quad\Longrightarrow\quad
 ds^2>0
\quad\text{(spacelike)}
}
\tag{A12}
\]

\[
\boxed{
\alpha=1
\quad\Longrightarrow\quad
 ds^2=0
\quad\text{(null)}
}
\tag{A13}
\]

\[
\boxed{
\alpha>1
\quad\Longrightarrow\quad
 ds^2<0
\quad\text{(timelike)}
}
\tag{A14}
\]

となる。

よって前報の特殊解は、空間尺度と周期尺度の指数が一致する

\[
\boxed{
\alpha=1
}
\]

という自己双対点であり、その点で全殻対区間が null となる。

また式 (A8) より、同じ規格化では

\[
\boxed{
\lambda_{jk}\nu_{jk}
=
\frac{1}{\alpha}
}
\tag{A15}
\]

である。次元を持つ基準速度 \[c_*\] を入れる規格化では右辺は
\[c_*/\alpha\] となる。

重要なのは、\[\lambda\nu\] 一定という性質そのものは null
に限定されないことである。全殻対で一定となることは冪関係
\[P\propto L^\alpha\] と同値であり、その中の \[\alpha=1\] が null
自己双対点を選ぶ。

\begin{center}\rule{0.5\linewidth}{0.5pt}\end{center}

\subsection{A.4　整数性と厳密閉包を保つ有理指数殻族}\label{a.4-ux6574ux6570ux6027ux3068ux53b3ux5bc6ux9589ux5305ux3092ux4fddux3064ux6709ux7406ux6307ux6570ux6bbbux65cf}

上の \[\alpha\]
一般化を、丸め誤差なしの可算複素ゼロ閉包解として構成する。

基礎整数列を

\[
B_k=k(k+1),
\qquad
k=1,2,3,\ldots
\tag{A16}
\]

とし、正整数 \[p,q\] に対して

\[
R_k=B_k^q,
\qquad
N_k=B_k^p
\tag{A17}
\]

を定義する。

各殻の複素波を

\[
\boxed{
z_{k,n}
=
\frac{A}{R_k}
\exp\left[
i\left(
\phi_k+
\frac{\pi n}{N_k}
\right)
\right],
\qquad
n=0,\ldots,N_k-1
}
\tag{A18}
\]

とする。

二乗すると

\[
z_{k,n}^2
=
\frac{A^2e^{2i\phi_k}}{R_k^2}
\exp\left(
\frac{2\pi i n}{N_k}
\right)
\]

なので、\[N_k\ge2\] なら単位根和により各殻は厳密に

\[
\boxed{
\sum_{n=0}^{N_k-1}z_{k,n}^2=0
}
\tag{A19}
\]

を満たす。

一方、振幅逆尺度から \[R_k\]、位相巡回次数または殻多重度から \[N_k\]
が読める。式 (A17) より

\[
\ln N_k
=
\frac{p}{q}
\ln R_k
\tag{A20}
\]

が厳密に成立する。

したがって

\[
\boxed{
\alpha=\frac{p}{q}
}
\tag{A21}
\]

という有理指数族が得られる。

前報の \[\alpha=1\] は \[p=q\] に対応する。

\subsubsection{絶対収束条件}\label{ux7d76ux5bfeux53ceux675fux6761ux4ef6}

全殻の二乗絶対値和は

\[
\sum_k\sum_{n=0}^{N_k-1}|z_{k,n}^2|
=
|A|^2
\sum_k
\frac{N_k}{R_k^2}
\]

であり、式 (A17) を代入すると

\[
|A|^2
\sum_k
B_k^{p-2q}
\tag{A22}
\]

となる。

\(B_k\sim k^2\) なので収束条件は

\[
2(p-2q)<-1.
\]

\(p,q\) は整数だから、これは

\[
\boxed{
p\le2q-1
}
\tag{A23}
\]

と同値である。

したがって絶対収束する有理指数族では

\[
\boxed{
0<\alpha
=
\frac pq
\le
2-
\frac1q
<2
}
\tag{A24}
\]

となる。

この範囲には、

\begin{itemize}
\tightlist
\item
  \(p<q\)：spacelike、
\item
  \(p=q\)：null、
\item
  \(q<p\le2q-1\)：timelike、
\end{itemize}

が全て含まれる。特に timelike な厳密族を持つには \[q\ge2\]
が必要である。

従って、前報の null
解を壊すことなく、整数巡回数・殻ごとの厳密二乗ゼロ閉包・全体絶対収束を同時に維持したまま、spacelike
と timelike の双方へ拡張できる。

\begin{center}\rule{0.5\linewidth}{0.5pt}\end{center}

\subsection{A.5　凝縮体との接続――導出済み事項と仮説の分離}\label{a.5-ux51ddux7e2eux4f53ux3068ux306eux63a5ux7d9aux5c0eux51faux6e08ux307fux4e8bux9805ux3068ux4eeeux8aacux306eux5206ux96e2}

本稿本体では、局在凝縮体を均質真空状態場の局所再配置として定義し、凝縮量、相対重心、有限広がり、および背景状態ゲージ歪みを
Fourier--Mellin 型読出しから構成した。

本補遺で新たに導出したのは次の事項である。

\begin{enumerate}
\def\labelenumi{\arabic{enumi}.}
\tightlist
\item
  全殻対で \[\lambda\nu\] が一定であることは、\[P\propto L^\alpha\]
  と必要十分に同値である。
\item
  \(\alpha=1\) は全殻対 null となる自己双対点である。
\item
  \(\alpha<1\) と \[\alpha>1\] は、それぞれ spacelike と timelike
  の読出しを与える。
\item
  有理指数 \[\alpha=p/q\]
  について、厳密二乗ゼロ閉包と絶対収束を保つ可算殻族が存在する。
\end{enumerate}

一方、次の命題は本稿ではまだ導出していない。

\begin{quote}
\textbf{接続仮説}：局在凝縮体の形成が、均質真空の \[\alpha=1\]
自己双対関係を局所的に破り、凝縮体内部またはその周囲に
\[\alpha_{\mathrm{eff}}>1\] の timelike 読出しを生じさせる可能性がある。
\end{quote}

局所的な有効指数は、二つの近接した読出し尺度に対して形式的に

\[
\boxed{
\alpha_{\mathrm{eff}}
=
\frac{\Delta\ln P}{\Delta\ln L}
}
\tag{A25}
\]

と定義できる。

均質真空で

\[
\alpha_{\mathrm{eff}}=1
\]

が保たれ、凝縮中心近傍で系統的に

\[
\alpha_{\mathrm{eff}}>1
\]

へ偏るなら、質量的凝縮と timelike 読出しの間に直接の関係が示唆される。

しかし現段階では、

\[
\boxed{
\text{局在凝縮}
\Longrightarrow
\alpha_{\mathrm{eff}}>1
}
\]

は仮説であり、帰結として扱わない。

さらに、\[L\] と \[P\] の役割交換が形式的に
\[\alpha\leftrightarrow1/\alpha\] を与えることから、自己双対点
\[\alpha=1\] を挟んで timelike 型と spacelike
型の読出しが逆数対を形成する可能性がある。これは標準的な massive-wave
に現れる位相速度と群速度の逆数構造を想起させるが、本稿では両者の物理的同一性を導出していない。従って、これは\textbf{対応仮説}に留める。

この仮説を検証する場合、\[\alpha_{\mathrm{eff}}\]
を「凝縮体の世界線型読出し」から測っているのか、「凝縮体に付随する波面型読出し」から測っているのかを区別しなければならない。両者を同じ
\[\alpha_{\mathrm{eff}}\] と無条件に同一視しない。

\begin{center}\rule{0.5\linewidth}{0.5pt}\end{center}

\subsection{A.6　最小判別計算}\label{a.6-ux6700ux5c0fux5224ux5225ux8a08ux7b97}

この接続仮説は、次の順序で判別できる。

\subsubsection{A.6.1　有理指数族の直接検算}\label{a.6.1-ux6709ux7406ux6307ux6570ux65cfux306eux76f4ux63a5ux691cux7b97}

複数の \[p,q\] について式 (A18) を生成し、

\begin{itemize}
\tightlist
\item
  各殻の二乗和が数値精度内でゼロ、
\item
  \(\ln N_k/\ln R_k=p/q\)、
\item
  \(\lambda\nu\) が全殻対で一定、
\item
  式 (A11) の符号が \[p/q\] の大小関係と一致、
\end{itemize}

することを確認する。

\subsubsection{A.6.2　局所閉包残差に対する読出し安定性}\label{a.6.2-ux5c40ux6240ux9589ux5305ux6b8bux5deeux306bux5bfeux3059ux308bux8aadux51faux3057ux5b89ux5b9aux6027}

二つの殻 \[j,k\] に複数の微小位相自由度を与え、

\[
Q_j=\epsilon,
\qquad
Q_k=-\epsilon,
\qquad
Q_{\mathrm{total}}=0
\tag{A26}
\]

を厳密に構成する。

一要素だけの位相変形

\[
\delta z
\simeq
iz\,\delta\theta
\]

では

\[
\delta(z^2)
\simeq
2iz^2\,\delta\theta
\tag{A27}
\]

となり、残差の複素方向が \[z^2\] に拘束されるため、任意の複素
\[\epsilon\]
を作るには一般に自由度が不足する。従って各殻で少なくとも二つの独立な実位相自由度を用意し、実部・虚部の二条件を解く。

その上で Fourier--Mellin 読出しの

\begin{itemize}
\tightlist
\item
  凝縮ピーク位置、
\item
  相対重心、
\item
  有限広がり、
\item
  \(\alpha_{\mathrm{eff}}\)、
\end{itemize}

が \[\epsilon\to0\]
に対して連続的に安定するかを測る。\[\alpha_{\mathrm{eff}}\]
については、世界線型読出しと波面型読出しのどちらから算出した値かを必ず記録し、両者が逆数関係を持つかどうかも独立に判別する。

この判別計算が安定なら、厳密特殊解が測度ゼロであっても、その近傍に物理的に読出し可能な構造が存在する可能性が高まる。

\begin{center}\rule{0.5\linewidth}{0.5pt}\end{center}

\subsection{A.7　補遺の結論}\label{a.7-ux88dcux907aux306eux7d50ux8ad6}

前報 {[}K1{]} の \[\alpha=1\] 特殊解は、単に \[\lambda\nu\]
が一定となる一例ではなく、冪関係

\[
P\propto L^\alpha
\]

の自己双対点であり、そこで全殻対が null となる。

本補遺では、

\[
\boxed{
\lambda\nu=\mathrm{const}
\quad\Longleftrightarrow\quad
P=P_0L^\alpha
}
\]

という必要十分条件を示し、さらに厳密な有理指数殻族により

\[
\boxed{
\alpha<1:
\text{spacelike},
\qquad
\alpha=1:
\text{null},
\qquad
\alpha>1:
\text{timelike}
}
\]

を同じ二乗ゼロ閉包系の中で構成した。

従って、null
構造は二乗ゼロ閉包そのものによって強制されるのではなく、空間尺度と周期尺度の自己双対条件
\[\alpha=1\] によって選ばれる。

局在凝縮がこの自己双対点からの局所偏差を生み、その偏差が質量的計量および
timelike
読出しと結び付くかどうかが、本稿から直接得られる次の判別課題である。

\begin{center}\rule{0.5\linewidth}{0.5pt}\end{center}

\subsection{参考文献}\label{ux53c2ux8003ux6587ux732e}

\subsubsection{自己引用}\label{ux81eaux5df1ux5f15ux7528}

{[}K1{]} Noriaki Kihara, \textbf{``Emergent Spacetime and Wave-Scale
Readout from Anonymous Complex Zero Closure --- Minimal Readout, an
Infinite Constructive Solution, Logarithmic Spacetime Mapping, and
Post-Mapping λν Invariance''}, 2026, v2.0. Japanese title:
\emph{無名複素ゼロ閉包からの時空・波動計量の読出し――最小読出し、無限閉包特殊解、対数時空写像、および写像後の
\[\lambda\nu\] 不変性}. Zenodo Concept DOI: 10.5281/zenodo.22282217.

\subsubsection{外部引用}\label{ux5916ux90e8ux5f15ux7528}

{[}E1{]} D. Casasent and D. Psaltis, ``Position, rotation, and scale
invariant optical correlation,'' \emph{Applied Optics}, \textbf{15}(7),
1795--1799 (1976). DOI: 10.1364/AO.15.001795.

{[}E2{]} Eugene Lukacs, \emph{Characteristic Functions}, 2nd ed., Hafner
Publishing Company / Griffin, 1970.

\begin{center}\rule{0.5\linewidth}{0.5pt}\end{center}

\subsection{付記：本稿の独自範囲}\label{ux4ed8ux8a18ux672cux7a3fux306eux72ecux81eaux7bc4ux56f2}

文献 {[}E1{]} は Fourier--Mellin
型のスケール・回転不変解析の既知性を、{[}E2{]}
は特性関数およびキュムラント展開からモーメントを得る標準数学を参照するために用いた。

これらの文献は、本稿の

\[
\sum_n z_n^2=0
\]

を満たす無名複素真空状態場、そこからの局在凝縮体の定義、凝縮重心を質量的計量として読む構成、あるいは凝縮による真空状態ゲージ歪みを提案するものではない。

本稿の独自提案は、これらの既知数学を、前報 {[}K1{]}
の無名複素ゼロ閉包からの創発時空読出しに接続し、

\[
\boxed{
\text{匿名複素状態}
\rightarrow
\text{凝縮体}
\rightarrow
\text{相対重心}
\rightarrow
\text{質量的計量}
\rightarrow
\text{背景ゲージ歪み}
}
\]

という一つの読出し系列として構成した点にある。

\end{document}
